\section{Structures argumentatives et mise en page de graphes}
Khartabil et collègues comparent trois visualisations et conduisent une étude à 21 participants. Leur problème central articule aperçu de la structure, lecture du contenu et navigation \cite{ref09}. Nous proposons un graphe principal et un panneau de passages synchronisé, afin de rendre ce compromis observable dans les essais utilisateurs.

La référence de Sugiyama, Tagawa et Toda est conservée dans la bibliographie comme point de départ historique pour approfondir la disposition hiérarchique \cite{ref10}. Le prototype choisit une disposition en couches orientée des raisons vers les conclusions ; l'ordre du document constitue une navigation complémentaire.
% TODO: vérifier le texte intégral et les métadonnées de ref10 ; accès IEEE limité.

StoryFlow combine optimisation discrète et continue pour disposer des storylines \cite{ref11a}. iStoryline propose des outils de génération et de réglage de ces représentations \cite{ref11b}. Leur continuité visuelle inspire la lecture progressive, mais l'axe logique de notre carte ne doit pas être interprété comme une chronologie.

Hurter, Ersoy et Telea regroupent les arêtes à partir d'une estimation de densité et de son gradient \cite{ref12}. Pour notre aperçu réduit, le regroupement n'est pas activé : nous faisons l'hypothèse que la lecture individuelle des relations prime. Son intérêt pourrait être évalué lorsque le lecteur développe plusieurs sous-processus.
